\documentclass[12pt, letterpaper]{article}

% ── Paquetes ──
\usepackage[utf8]{inputenc}
\usepackage[spanish]{babel}
\usepackage[margin=2.5cm]{geometry}
\usepackage{graphicx}
\usepackage{listings}
\usepackage{xcolor}
\usepackage{hyperref}
\usepackage{fancyhdr}
\usepackage{amsmath}
\usepackage{float}
\usepackage{parskip}

% ── Estilo de código C ──
\definecolor{codegreen}{rgb}{0,0.6,0}
\definecolor{codegray}{rgb}{0.5,0.5,0.5}
\definecolor{codepurple}{rgb}{0.58,0,0.82}
\definecolor{backcolour}{rgb}{0.97,0.97,0.97}

\lstdefinestyle{cstyle}{
    backgroundcolor=\color{backcolour},
    commentstyle=\color{codegreen},
    keywordstyle=\color{blue}\bfseries,
    numberstyle=\tiny\color{codegray},
    stringstyle=\color{codepurple},
    basicstyle=\ttfamily\footnotesize,
    breakatwhitespace=false,
    breaklines=true,
    captionpos=b,
    keepspaces=true,
    numbers=left,
    numbersep=5pt,
    showspaces=false,
    showstringspaces=false,
    showtabs=false,
    tabsize=4,
    language=C,
    frame=single,
    rulecolor=\color{codegray}
}
\lstset{style=cstyle,
    literate=
        {á}{{\'a}}1 {é}{{\'e}}1 {í}{{\'i}}1 {ó}{{\'o}}1 {ú}{{\'u}}1
        {Á}{{\'A}}1 {É}{{\'E}}1 {Í}{{\'I}}1 {Ó}{{\'O}}1 {Ú}{{\'U}}1
        {ñ}{{\~n}}1 {Ñ}{{\~N}}1
        {π}{{$\pi$}}1
}

% ── Encabezado y pie de página ──
\setlength{\headheight}{14.5pt}
\addtolength{\topmargin}{-2.5pt}
\pagestyle{fancy}
\fancyhf{}
\fancyhead[L]{Taller 1 -- Programación Estructurada y Archivos en C}
\fancyhead[R]{2026-1}
\fancyfoot[C]{\thepage}

% ══════════════════════════════════════════════
\begin{document}

% ── Portada ──
\begin{titlepage}
    \centering
    \vspace*{2cm}

    {\Huge\bfseries Taller 1: Programación Estructurada y Archivos en C\par}
    \vspace{1.5cm}

    {\Large Introducción a la Programación Embebida\par}
    \vspace{2cm}

    {\large\textbf{Docente:} Carlos Andrés Erazo Garzón\\
    \href{mailto:carlos-erazo@javeriana.edu.co}{carlos-erazo@javeriana.edu.co}\par}
    \vspace{2cm}

    {\large\textbf{Estudiante:}\\[0.3cm]
    Juan Pablo Cañón Contreras\par}
    % Si hay más integrantes, agregar aquí:
    % Nombre Apellido Apellido\\
    % Nombre Apellido Apellido\par}
    \vspace{2cm}

    {\large Departamento de Electrónica\\
    Pontificia Universidad Javeriana\par}
    \vspace{1cm}

    {\large Semestre 2026-1\par}

    \vfill
    {\large \today\par}
\end{titlepage}

% ── Tabla de contenidos ──
\tableofcontents
\newpage

% ══════════════════════════════════════════════
\section{Introducción}

El presente informe documenta el desarrollo del Taller~1 de la asignatura
\textit{Introducción a la Programación Embebida}. El objetivo principal es
aplicar los conceptos de tipos de datos, cadenas de caracteres, arreglos (1D y
2D), estructuras (\texttt{struct}) y persistencia de datos (lectura y escritura
de archivos con \texttt{sscanf} y \texttt{fprintf}) mediante un diseño modular
basado en funciones.

Se sigue la metodología de programación estructurada en~C, donde cada punto del
taller se resuelve mediante funciones independientes que reciben parámetros
validados y retornan resultados, manteniendo el \texttt{main} exclusivamente
para la interacción con el usuario y la invocación de dichas funciones.

% ══════════════════════════════════════════════
\section{Punto 1: Aproximación de $\pi$ con la serie de Leibniz}

\subsection{Descripción del problema}

Aproximar el valor de $\pi$ mediante la serie de Leibniz:

\[
    \pi \approx 4 \sum_{k=0}^{N} \frac{(-1)^k}{2k+1}
    = 4 \left(1 - \frac{1}{3} + \frac{1}{5} - \frac{1}{7} + \cdots \right)
\]

El usuario ingresa el número de términos $N$ y el programa calcula la
aproximación correspondiente.

\subsection{Funciones implementadas}

\begin{table}[H]
\centering
\begin{tabular}{|l|l|l|p{6cm}|}
\hline
\textbf{Función} & \textbf{Entrada} & \textbf{Retorno} & \textbf{Descripción} \\
\hline
\texttt{calcular\_pi} & \texttt{int N} & \texttt{double} &
Calcula la aproximación de $\pi$ sumando $N+1$ términos de la serie de
Leibniz. Alterna suma y resta según la paridad de $k$. \\
\hline
\end{tabular}
\caption{Funciones del Punto~1}
\end{table}

\subsection{Validación de datos}

\begin{itemize}
    \item Se verifica que $N \geq 0$. Si el usuario ingresa un valor negativo,
          se imprime un mensaje de error y el programa termina de forma
          controlada.
\end{itemize}

\subsection{Código fuente}

\lstinputlisting[caption={punto1.c -- Aproximación de $\pi$}]{punto1.c}

\subsection{Casos de prueba}

\begin{table}[H]
\centering
\begin{tabular}{|c|c|c|}
\hline
\textbf{Entrada ($N$)} & \textbf{Salida (Aproximación de $\pi$)} & \textbf{Observación} \\
\hline
0   & 4.00000000 & Solo el primer término \\
1   & 2.66666667 & Dos términos \\
10  & 3.23231580 & Convergencia lenta \\
100 & 3.15149340 & Mejor aproximación \\
1000 & 3.14259165 & Cercano a $\pi$ \\
$-5$ & ERROR: N debe ser $\geq 0$ & Validación de entrada \\
\hline
\end{tabular}
\caption{Casos de prueba del Punto~1}
\end{table}

% ══════════════════════════════════════════════
\section{Punto 2: Cifrado César}

\subsection{Descripción del problema}

Implementar el cifrado César sobre una cadena de caracteres. El usuario ingresa
una frase y un valor de desplazamiento $d$, y el programa sustituye cada letra
por la que se encuentra $d$ posiciones más adelante en el alfabeto (con
envolvimiento circular). Los caracteres que no son letras (espacios, números,
signos) se mantienen sin cambio.

\[
    c' = \bigl((c - \text{base}) + d\bigr) \bmod 26 + \text{base}
\]

donde \texttt{base} es \texttt{'A'} o \texttt{'a'} según el caso.

\subsection{Funciones implementadas}

\begin{table}[H]
\centering
\begin{tabular}{|l|p{4cm}|l|p{5cm}|}
\hline
\textbf{Función} & \textbf{Entrada} & \textbf{Retorno} & \textbf{Descripción} \\
\hline
\texttt{cifrar\_cadena} &
\texttt{char texto[]},
\texttt{int desplazamiento} &
\texttt{void} &
Recorre la cadena carácter a carácter y aplica el desplazamiento César
solo a letras mayúsculas y minúsculas. Modifica el arreglo in-place. \\
\hline
\end{tabular}
\caption{Funciones del Punto~2}
\end{table}

\subsection{Validación de datos}

\begin{itemize}
    \item Se soportan desplazamientos negativos: si el resultado del módulo
          es negativo, se suma 26 para mantener el rango válido.
    \item Los caracteres no alfabéticos (espacios, números, signos) no se
          modifican.
    \item Se usa \texttt{fgets} para leer la frase (evita desbordamiento) y
          se elimina el salto de línea residual.
\end{itemize}

\subsection{Código fuente}

\lstinputlisting[caption={punto2.c -- Cifrado César}]{punto2.c}

\subsection{Casos de prueba}

\begin{table}[H]
\centering
\begin{tabular}{|l|c|l|l|}
\hline
\textbf{Frase} & \textbf{Despl.} & \textbf{Salida} & \textbf{Observación} \\
\hline
Hola Mundo & 3 & Krod Pxqgr & Desplazamiento positivo \\
\hline
Krod Pxqgr & $-3$ & Hola Mundo & Descifrado inverso \\
\hline
ABC xyz & 1 & BCD yza & Envolvimiento z$\to$a \\
\hline
Hola 123! & 5 & Mtqf 123! & No altera no-letras \\
\hline
\end{tabular}
\caption{Casos de prueba del Punto~2}
\end{table}

% ══════════════════════════════════════════════
\section{Punto 3: Filtro de promedio móvil (suavizado de señal)}

\subsection{Descripción del problema}

Implementar un filtro de promedio móvil que suavice un arreglo de datos
simulando una señal con ruido. Dado un arreglo de entrada y un tamaño de
ventana $W$, cada elemento de salida se calcula como el promedio de los
últimos $W$ elementos disponibles (o menos, si no hay suficientes hacia
la izquierda).

\[
    \text{salida}[i] = \frac{1}{|\text{ventana}|} \sum_{j=\text{inicio}}^{i} \text{entrada}[j]
\]

donde $\text{inicio} = \max(0,\; i - W + 1)$.

\subsection{Funciones implementadas}

\begin{table}[H]
\centering
\begin{tabular}{|l|p{4.5cm}|l|p{4.5cm}|}
\hline
\textbf{Función} & \textbf{Entrada} & \textbf{Retorno} & \textbf{Descripción} \\
\hline
\texttt{filtro\_promedio} &
\texttt{float entrada[]},
\texttt{float salida[]},
\texttt{int tamano},
\texttt{int ventana} &
\texttt{void} &
Aplica un filtro de promedio móvil sobre el arreglo \texttt{entrada} y
escribe el resultado suavizado en \texttt{salida}. \\
\hline
\end{tabular}
\caption{Funciones del Punto~3}
\end{table}

\subsection{Validación de datos}

\begin{itemize}
    \item Se verifica que \texttt{ventana} y \texttt{tamano} sean mayores que 0.
          Si alguno es inválido, se imprime un mensaje de error y la función
          retorna sin modificar el arreglo de salida.
\end{itemize}

\subsection{Código fuente}

\lstinputlisting[caption={punto3.c -- Filtro de promedio móvil}]{punto3.c}

\subsection{Casos de prueba}

\begin{table}[H]
\centering
\small
\begin{tabular}{|p{3.5cm}|c|p{3.5cm}|p{3.5cm}|}
\hline
\textbf{Entrada} & \textbf{W} & \textbf{Salida (suavizada)} & \textbf{Observación} \\
\hline
\{10.2, 10.5, 25.0, 10.3, 10.1\} & 3 &
\{10.2, 10.4, 15.2, 15.3, 15.1\} & El pico 25.0 se atenúa \\
\hline
\{10.2, 10.5, 25.0, 10.3, 10.1\} & 1 &
\{10.2, 10.5, 25.0, 10.3, 10.1\} & Sin suavizado \\
\hline
\{10.2, 10.5, 25.0, 10.3, 10.1\} & 5 &
\{10.2, 10.4, 15.2, 14.0, 13.2\} & Máximo suavizado \\
\hline
Cualquiera & 0 &
Error & Validación de ventana \\
\hline
\end{tabular}
\caption{Casos de prueba del Punto~3}
\end{table}

% ══════════════════════════════════════════════
\section{Punto 4: Detección de picos en una matriz 2D}

\subsection{Descripción del problema}

Dada una matriz de enteros de $4\times4$ que representa un mapa de valores
(por ejemplo, intensidades de señal), detectar todos los \textit{picos}: celdas
cuyo valor es estrictamente mayor que el de sus cuatro vecinos ortogonales
(arriba, abajo, izquierda, derecha). Las celdas en los bordes tienen menos
vecinos, lo cual se maneja verificando los límites antes de comparar.

\subsection{Funciones implementadas}

\begin{table}[H]
\centering
\begin{tabular}{|l|p{4cm}|l|p{5cm}|}
\hline
\textbf{Función} & \textbf{Entrada} & \textbf{Retorno} & \textbf{Descripción} \\
\hline
\texttt{detectar\_picos} &
\texttt{int filas},
\texttt{int cols},
\texttt{int mapa[4][4]} &
\texttt{void} &
Recorre la matriz comparando cada celda con sus vecinos. Almacena las
coordenadas de los picos encontrados y los imprime. \\
\hline
\end{tabular}
\caption{Funciones del Punto~4}
\end{table}

\subsection{Validación de datos}

\begin{itemize}
    \item Antes de acceder a cada vecino se verifica que el índice esté
          dentro de los límites de la matriz (\texttt{i-1 >= 0},
          \texttt{j+1 < cols}, etc.).
    \item Si no se encuentran picos, se imprime un mensaje informativo.
\end{itemize}

\subsection{Código fuente}

\lstinputlisting[caption={punto4.c -- Detección de picos en matriz 2D}]{punto4.c}

\subsection{Casos de prueba}

\begin{table}[H]
\centering
\begin{tabular}{|p{5cm}|p{5cm}|p{3.5cm}|}
\hline
\textbf{Matriz de entrada} & \textbf{Salida} & \textbf{Observación} \\
\hline
\texttt{\{1,2,1,0\}, \{2,9,2,0\},} \texttt{\{1,2,1,0\}, \{0,0,0,0\}} &
Pico en (1,1) valor 9 &
Mayor que sus 4 vecinos \\
\hline
Matriz con todos los valores iguales &
No se encontraron picos &
Ningún valor es estrictamente mayor \\
\hline
\end{tabular}
\caption{Casos de prueba del Punto~4}
\end{table}

% ══════════════════════════════════════════════
\section{Punto 5: Lectura de archivo con estructuras}

\subsection{Descripción del problema}

Leer un archivo de texto (\texttt{sensores.txt}) que contiene mediciones de
sensores en formato CSV (\texttt{ID,temperatura,humedad}) y almacenar cada
registro en un arreglo de estructuras \texttt{Medicion}. Luego mostrar los
datos cargados por pantalla.

Formato del archivo:
\begin{verbatim}
SENS_01,24.5,60
SENS_02,22.3,55
SENS_03,25.0,62
\end{verbatim}

\subsection{Funciones implementadas}

\begin{table}[H]
\centering
\begin{tabular}{|l|p{4cm}|l|p{5cm}|}
\hline
\textbf{Función} & \textbf{Entrada} & \textbf{Retorno} & \textbf{Descripción} \\
\hline
\texttt{cargar\_datos} &
\texttt{char nombre\_archivo[]},
\texttt{Medicion datos[]},
\texttt{int max\_datos} &
\texttt{int} &
Abre el archivo, lee línea a línea con \texttt{fgets} y extrae los campos
con \texttt{sscanf}. Retorna la cantidad de registros leídos, o $-1$ si
el archivo no se pudo abrir. \\
\hline
\end{tabular}
\caption{Funciones del Punto~5}
\end{table}

\subsection{Estructura de datos}

\begin{verbatim}
typedef struct {
    char id[32];
    float temperatura;
    int humedad;
} Medicion;
\end{verbatim}

\subsection{Validación de datos}

\begin{itemize}
    \item Si \texttt{fopen} retorna \texttt{NULL}, se imprime un mensaje de
          error y se retorna $-1$ sin colapsar.
    \item No se leen más de \texttt{max\_datos} registros para evitar
          desbordamiento del arreglo.
    \item Solo se guardan líneas donde \texttt{sscanf} parsea exactamente
          3 campos.
\end{itemize}

\subsection{Código fuente}

\lstinputlisting[caption={punto5.c -- Lectura de archivo con estructuras}]{punto5.c}

\subsection{Casos de prueba}

\begin{table}[H]
\centering
\small
\begin{tabular}{|p{3.5cm}|p{5.5cm}|p{4cm}|}
\hline
\textbf{Archivo} & \textbf{Salida} & \textbf{Observación} \\
\hline
\texttt{sensores.txt} (3 líneas válidas) &
Se cargaron 3 mediciones: SENS\_01 24.5 60, SENS\_02 22.3 55, SENS\_03 25.0 62 &
Lectura correcta \\
\hline
Archivo inexistente &
Error: no se pudo abrir el archivo &
Manejo de \texttt{fopen} NULL \\
\hline
\end{tabular}
\caption{Casos de prueba del Punto~5}
\end{table}

% ══════════════════════════════════════════════
\section{Punto 6: Análisis de temperaturas}

\subsection{Descripción del problema}

A partir de los datos cargados desde \texttt{sensores.txt} (reutilizando
\texttt{cargar\_datos} del Punto~5), determinar la temperatura máxima, la
mínima y el promedio general, indicando el ID del sensor correspondiente para
los extremos.

\subsection{Funciones implementadas}

\begin{table}[H]
\centering
\begin{tabular}{|l|p{3.5cm}|l|p{5cm}|}
\hline
\textbf{Función} & \textbf{Entrada} & \textbf{Retorno} & \textbf{Descripción} \\
\hline
\texttt{analizar\_temperaturas} &
\texttt{Medicion datos[]},
\texttt{int cantidad} &
\texttt{void} &
Recorre el arreglo de mediciones para encontrar la temperatura máxima,
mínima y el promedio. Imprime los resultados con el ID del sensor. \\
\hline
\texttt{cargar\_datos} &
(igual que Punto~5) &
\texttt{int} &
Reutilizada del Punto~5. \\
\hline
\end{tabular}
\caption{Funciones del Punto~6}
\end{table}

\subsection{Validación de datos}

\begin{itemize}
    \item Si \texttt{cantidad} es $\leq 0$, se imprime ``No hay datos para
          analizar'' y la función retorna sin operar (evita división por cero
          en el promedio).
\end{itemize}

\subsection{Código fuente}

\lstinputlisting[caption={punto6.c -- Análisis de temperaturas}]{punto6.c}

\subsection{Casos de prueba}

\begin{table}[H]
\centering
\begin{tabular}{|p{3.5cm}|p{6cm}|p{3.5cm}|}
\hline
\textbf{Archivo} & \textbf{Salida} & \textbf{Observación} \\
\hline
\texttt{sensores.txt} (3 sensores) &
Máxima: 25.0 (SENS\_03), Mínima: 22.3 (SENS\_02), Promedio: 23.9 &
Cálculos correctos \\
\hline
Archivo vacío o inexistente &
Error o ``No hay datos'' &
Validación de entrada \\
\hline
\end{tabular}
\caption{Casos de prueba del Punto~6}
\end{table}

% ══════════════════════════════════════════════
\section{Punto 7: Exportar datos corregidos a archivo}

\subsection{Descripción del problema}

Cargar las mediciones desde \texttt{sensores.txt}, aplicar un \textit{offset}
de calibración a todas las temperaturas y escribir los datos corregidos en un
nuevo archivo (\texttt{sensores\_calibrados.txt}) con el mismo formato CSV.

\subsection{Funciones implementadas}

\begin{table}[H]
\centering
\begin{tabular}{|l|p{3.5cm}|l|p{5cm}|}
\hline
\textbf{Función} & \textbf{Entrada} & \textbf{Retorno} & \textbf{Descripción} \\
\hline
\texttt{exportar\_corregido} &
\texttt{char nom\_arch[]},
\texttt{Medicion datos[]},
\texttt{int cant},
\texttt{float offset} &
\texttt{void} &
Crea un archivo de salida y escribe cada medición con la temperatura
corregida (\texttt{temp + offset}) usando \texttt{fprintf}. \\
\hline
\texttt{cargar\_datos} &
(igual que Punto~5) &
\texttt{int} &
Reutilizada del Punto~5. \\
\hline
\end{tabular}
\caption{Funciones del Punto~7}
\end{table}

\subsection{Validación de datos}

\begin{itemize}
    \item Si \texttt{fopen} para escritura retorna \texttt{NULL}, se imprime
          un mensaje de error y la función retorna sin colapsar.
    \item Se verifica el retorno de \texttt{cargar\_datos} antes de proceder.
\end{itemize}

\subsection{Código fuente}

\lstinputlisting[caption={punto7.c -- Exportar datos corregidos}]{punto7.c}

\subsection{Casos de prueba}

\begin{table}[H]
\centering
\small
\begin{tabular}{|c|p{2.8cm}|p{2.8cm}|p{3.8cm}|}
\hline
\textbf{Offset} & \textbf{Temp. original} & \textbf{Temp. corregida} & \textbf{Observación} \\
\hline
$-1.5$ & 24.5, 22.3, 25.0 & 23.0, 20.8, 23.5 &
Genera archivo calibrado \\
\hline
$+2.0$ & 24.5, 22.3, 25.0 & 26.5, 24.3, 27.0 &
Offset positivo \\
\hline
$0.0$ & 24.5, 22.3, 25.0 & 24.5, 22.3, 25.0 &
Sin cambio \\
\hline
\end{tabular}
\caption{Casos de prueba del Punto~7}
\end{table}

% ══════════════════════════════════════════════
\section{Conclusiones}

\begin{itemize}
    \item Se aplicó el diseño modular separando la lógica en funciones
          independientes, manteniendo el \texttt{main} limpio y enfocado en la
          interacción con el usuario.
    \item Se implementó validación de datos en todas las entradas del usuario
          y en los parámetros de las funciones, garantizando que el programa
          no colapse ante valores inválidos.
    \item Se manipularon cadenas de caracteres mediante aritmética de posiciones
          ASCII para implementar el cifrado César, con soporte para
          desplazamientos negativos.
    \item Se trabajó con arreglos 1D para el filtro de promedio móvil y con
          arreglos 2D para la detección de picos, aplicando verificación
          de límites en los accesos a vecinos.
    \item Se utilizaron estructuras (\texttt{struct}) para organizar datos
          de sensores de forma coherente, agrupando ID, temperatura y humedad
          en un solo tipo de dato.
    \item Se practicó la persistencia de datos mediante lectura y escritura
          de archivos con \texttt{fopen}, \texttt{fgets}, \texttt{sscanf} y
          \texttt{fprintf}, verificando siempre el retorno de \texttt{fopen}
          para evitar errores fatales.
    \item Se demostró la reutilización de funciones al emplear
          \texttt{cargar\_datos} en los puntos 5, 6 y 7, evidenciando las
          ventajas del diseño modular.
\end{itemize}

\end{document}